% GENERATED FILE. Edit canonical lessons, then run npm run generate:books.
% canonical-source-hash: fnv1a64:92a30b8f694f4a2e
% canonical-lessons: HI-C444-bandargah, HI-C444-langar, HI-C444-khari, HI-C444-tat, HI-C444-dubki, HI-R444-review-438-444

\chapter{Harbours and Bays}
\label{ch:hi-harbours-and-bays}
\hlchaptermodality{444}

\begin{chapteropening}
\textbf{By the end of this chapter:} \emph{I can talk about a harbour, an anchor, a bay and a coast, and a cold dip in the sea.}

Complete the last lesson of chapter 444: I can talk about a harbour, an anchor, a bay and a coast, and a cold dip in the sea.
\end{chapteropening}

\section[\texorpdfstring{\hi{बंदरगाह} (bandargāh)}{बंदरगाह (bandargāh)}]{\hi{बंदरगाह} (bandargāh) — a harbour, a port}

\label{lesson:HI-C444-bandargah}



\begin{quote}
Before the new one: say the Hindi for to irrigate, then the Hindi for to weed.
\end{quote}

\subsection*{You'll want to know: \hi{बंदरगाह}}
\textbf{\hi{बंदरगाह}} — \emph{bandargāh} — \textquotedblleft{}a harbour, a port\textquotedblright{}.

\begin{quote}
\hi{बंदरगाह} \hi{पर} \hi{बड़े} \hi{जहाज़} \hi{खड़े} \hi{थे।}
\end{quote}

\begin{quote}
\emph{bandargāh par baṛe jahāz khaṛe the.}
\end{quote}

\begin{quote}
Big ships were standing in the harbour.
\end{quote}

Masculine.

\begin{cousinweb}
\textbf{\hi{बंदरगाह}} (\emph{bandargāh}) is usually traced to Persian \emph{bandargāh}, from \emph{bandar}, \textquotedblleft{}a port\textquotedblright{}, and \emph{-gāh}, \textquotedblleft{}a place\textquotedblright{}.
\end{cousinweb}

\begin{grammarlens}[title={What you've built}]
One more word for getting around.
\end{grammarlens}

\subsection*{Your turn}
\begin{itemize}
\raggedright
  \item \emph{Say it:} \emph{bandargāh}
  \item \emph{Say it:} \emph{bandargāh}, once more
  \item \emph{Say it:} \emph{bandargāh par baṛe jahāz khaṛe the.} — Big ships were standing in the harbour.
  \item \emph{Recall:} say the Hindi for to irrigate, then the Hindi for to weed, then say \emph{bandargāh} again
\end{itemize}

\begin{tcolorbox}[breakable,colback=teal!4,colframe=teal!35!black,title={Before you move on}]
What does \hi{बंदरगाह} mean? (\textquotedblleft{}A harbour, a port\textquotedblright{}.) Say it once more.
\end{tcolorbox}
\section[\texorpdfstring{\hi{लंगर} (laṅgar)}{लंगर (laṅgar)}]{\hi{लंगर} (laṅgar) — an anchor; a free kitchen}

\label{lesson:HI-C444-langar}



\begin{quote}
Before the new one: say the Hindi for to weed, then the Hindi for a harbour.
\end{quote}

\subsection*{You'll want to know: \hi{लंगर}}
\textbf{\hi{लंगर}} — \emph{laṅgar} — \textquotedblleft{}an anchor; a free kitchen\textquotedblright{}.

\begin{quote}
\hi{जहाज़} \hi{का} \hi{लंगर} \hi{बहुत} \hi{भारी} \hi{था।}
\end{quote}

\begin{quote}
\emph{jahāz kā laṅgar bahut bhārī thā.}
\end{quote}

\begin{quote}
The ship's anchor was very heavy.
\end{quote}

Masculine.

At a gurdwara, the \hi{लंगर} is the free kitchen where everyone eats together.

\begin{cousinweb}
\textbf{\hi{लंगर}} (\emph{laṅgar}) is usually traced to Persian \emph{langar}, \textquotedblleft{}an anchor\textquotedblright{}, and also a place where food is given out; both senses live in Hindi.
\end{cousinweb}

\begin{grammarlens}[title={What you've built}]
One more word for getting around.
\end{grammarlens}

\subsection*{Your turn}
\begin{itemize}
\raggedright
  \item \emph{Say it:} \emph{laṅgar}
  \item \emph{Say it:} \emph{laṅgar}, once more
  \item \emph{Say it:} \emph{jahāz kā laṅgar bahut bhārī thā.} — The ship's anchor was very heavy.
  \item \emph{Recall:} say the Hindi for to weed, then the Hindi for a harbour, then say \emph{laṅgar} again
\end{itemize}

\begin{tcolorbox}[breakable,colback=teal!4,colframe=teal!35!black,title={Before you move on}]
What does \hi{लंगर} mean? (\textquotedblleft{}An anchor; a free kitchen\textquotedblright{}.) Say it once more.
\end{tcolorbox}
\section[\texorpdfstring{\hi{खाड़ी} (khāṛī)}{खाड़ी (khāṛī)}]{\hi{खाड़ी} (khāṛī) — a bay, a gulf}

\label{lesson:HI-C444-khari}



\begin{quote}
Before the new one: say the Hindi for a harbour, then the Hindi for an anchor.
\end{quote}

\subsection*{You'll want to know: \hi{खाड़ी}}
\textbf{\hi{खाड़ी}} — \emph{khāṛī} — \textquotedblleft{}a bay, a gulf\textquotedblright{}.

\begin{quote}
\hi{हमारा} \hi{गाँव} \hi{खाड़ी} \hi{के} \hi{किनारे} \hi{है।}
\end{quote}

\begin{quote}
\emph{hamārā gā̃v khāṛī ke kināre hai.}
\end{quote}

\begin{quote}
Our village is on the shore of the bay.
\end{quote}

Feminine.

\begin{cousinweb}
\textbf{\hi{खाड़ी}} (\emph{khāṛī}) goes back to \emph{*khaḍḍa}, \textquotedblleft{}a hole, a pit\textquotedblright{}, a form CDIAL reconstructs (CDIAL 3790), which lists Hindi \emph{khāṛī}, \textquotedblleft{}a creek, an inlet\textquotedblright{}: the hollow the sea fills.
\end{cousinweb}

\begin{grammarlens}[title={What you've built}]
One more word for getting around.
\end{grammarlens}

\subsection*{Your turn}
\begin{itemize}
\raggedright
  \item \emph{Say it:} \emph{khāṛī}
  \item \emph{Say it:} \emph{khāṛī}, once more
  \item \emph{Say it:} \emph{hamārā gā̃v khāṛī ke kināre hai.} — Our village is on the shore of the bay.
  \item \emph{Recall:} say the Hindi for a harbour, then the Hindi for an anchor, then say \emph{khāṛī} again
\end{itemize}

\begin{tcolorbox}[breakable,colback=teal!4,colframe=teal!35!black,title={Before you move on}]
What does \hi{खाड़ी} mean? (\textquotedblleft{}A bay, a gulf\textquotedblright{}.) Say it once more.
\end{tcolorbox}
\section[\texorpdfstring{\hi{तट} (taṭ)}{तट (taṭ)}]{\hi{तट} (taṭ) — a coast; a bank}

\label{lesson:HI-C444-tat}



\begin{quote}
Before the new one: say the Hindi for an anchor, then the Hindi for a bay.
\end{quote}

\subsection*{You'll want to know: \hi{तट}}
\textbf{\hi{तट}} — \emph{taṭ} — \textquotedblleft{}a coast; a bank\textquotedblright{}.

\begin{quote}
\hi{हम} \hi{समुद्र} \hi{के} \hi{तट} \hi{पर} \hi{टहले।}
\end{quote}

\begin{quote}
\emph{ham samudra ke taṭ par ṭahle.}
\end{quote}

\begin{quote}
We strolled along the seashore.
\end{quote}

Masculine.

\begin{cousinweb}
\textbf{\hi{तट}} (\emph{taṭ}) is a tatsama, Sanskrit \emph{taṭa}, \textquotedblleft{}a slope, a bank, a shore\textquotedblright{} (Monier-Williams): the formal partner of \hi{किनारा}.
\end{cousinweb}

\begin{grammarlens}[title={What you've built}]
One more word for getting around.
\end{grammarlens}

\subsection*{Your turn}
\begin{itemize}
\raggedright
  \item \emph{Say it:} \emph{taṭ}
  \item \emph{Say it:} \emph{taṭ}, once more
  \item \emph{Say it:} \emph{ham samudra ke taṭ par ṭahle.} — We strolled along the seashore.
  \item \emph{Recall:} say the Hindi for an anchor, then the Hindi for a bay, then say \emph{taṭ} again
\end{itemize}

\begin{tcolorbox}[breakable,colback=teal!4,colframe=teal!35!black,title={Before you move on}]
What does \hi{तट} mean? (\textquotedblleft{}A coast; a bank\textquotedblright{}.) Say it once more.
\end{tcolorbox}
\section[\texorpdfstring{\hi{डुबकी} (ḍubkī)}{डुबकी (ḍubkī)}]{\hi{डुबकी} (ḍubkī) — a dip (in water)}

\label{lesson:HI-C444-dubki}



\begin{quote}
Before the new one: say the Hindi for a bay, then the Hindi for a coast.
\end{quote}

\subsection*{You'll want to know: \hi{डुबकी}}
\textbf{\hi{डुबकी}} — \emph{ḍubkī} — \textquotedblleft{}a dip (in water)\textquotedblright{}.

\begin{quote}
\hi{ठंडे} \hi{पानी} \hi{में} \hi{एक} \hi{डुबकी} \hi{के} \hi{बाद} \hi{मैं} \hi{काँप} \hi{रहा} \hi{था।}
\end{quote}

\begin{quote}
\emph{ṭhaṇḍe pānī meṁ ek ḍubkī ke bād maiṁ kā̃p rahā thā.}
\end{quote}

\begin{quote}
After one dip in the cold water I was shivering.
\end{quote}

Feminine.

\begin{cousinweb}
\textbf{\hi{डुबकी}} (\emph{ḍubkī}) is built on the stem of \hi{डूबना}, \textquotedblleft{}to sink\textquotedblright{}, from \emph{*ḍubb}, \textquotedblleft{}sink\textquotedblright{}, a form CDIAL reconstructs (CDIAL 5561); CDIAL does not list the noun itself.
\end{cousinweb}

\begin{grammarlens}[title={What you've built}]
One more word for getting around.
\end{grammarlens}

\subsection*{Your turn}
\begin{itemize}
\raggedright
  \item \emph{Say it:} \emph{ḍubkī}
  \item \emph{Say it:} \emph{ḍubkī}, once more
  \item \emph{Say it:} \emph{ṭhaṇḍe pānī meṁ ek ḍubkī ke bād maiṁ kā̃p rahā thā.} — After one dip in the cold water I was shivering.
  \item \emph{Recall:} say the Hindi for a bay, then the Hindi for a coast, then say \emph{ḍubkī} again
\end{itemize}

\begin{tcolorbox}[breakable,colback=teal!4,colframe=teal!35!black,title={Before you move on}]
What does \hi{डुबकी} mean? (\textquotedblleft{}A dip (in water)\textquotedblright{}.) Say it once more.
\end{tcolorbox}
\section[(dialogue)]{Looking back: from Winter Cold to Harbours and Bays}

\label{lesson:HI-R444-review-438-444}



\begin{quote}
Say the Hindi for winter and a hailstone.
\end{quote}

\subsection*{Your turn}
Say these aloud:

\begin{itemize}
\raggedright
  \item winter, a hailstone, cold (formal), snowfall, to wrap oneself in
  \item heat (formal), mercury, temperature, a degree, to scorch
  \item dust, to be surrounded, to let up, to clear (of clouds), to blow away
  \item daybreak, sunrise, sunset, the horizon, dusk
  \item a summit, a mound, a heap, a straw, to roll (something)
  \item a canal, a reservoir, a current, to irrigate, to weed
  \item a harbour, an anchor, a bay, a coast, a dip
\end{itemize}

\begin{tcolorbox}[breakable,colback=teal!4,colframe=teal!35!black,title={Before you move on}]
Winter? (\textbf{\hi{जाड़ा}}, \emph{jāṛā}.) A hailstone? (\textbf{\hi{ओला}}, \emph{olā}.) Cold (formal)? (\textbf{\hi{शीत}}, \emph{śīt}.) Snowfall? (\textbf{\hi{हिमपात}}, \emph{himpāt}.) To wrap oneself in? (\textbf{\hi{ओढ़ना}}, \emph{oṛhnā}.) Heat (formal)? (\textbf{\hi{ताप}}, \emph{tāp}.) Mercury? (\textbf{\hi{पारा}}, \emph{pārā}.) Temperature? (\textbf{\hi{तापमान}}, \emph{tāpmān}.) A degree? (\textbf{\hi{डिग्री}}, \emph{ḍigrī}.) To scorch? (\textbf{\hi{तपना}}, \emph{tapnā}.) Dust? (\textbf{\hi{धूल}}, \emph{dhūl}.) To be surrounded? (\textbf{\hi{घिरना}}, \emph{ghirnā}.) To let up? (\textbf{\hi{थमना}}, \emph{thamnā}.) To clear (of clouds)? (\textbf{\hi{छँटना}}, \emph{chãṭnā}.) To blow away? (\textbf{\hi{उड़ाना}}, \emph{uṛānā}.) Daybreak? (\textbf{\hi{सवेरा}}, \emph{saverā}.) Sunrise? (\textbf{\hi{सूर्योदय}}, \emph{sūryoday}.) Sunset? (\textbf{\hi{सूर्यास्त}}, \emph{sūryāst}.) The horizon? (\textbf{\hi{क्षितिज}}, \emph{kṣitij}.) Dusk? (\textbf{\hi{गोधूलि}}, \emph{godhūli}.) A summit? (\textbf{\hi{शिखर}}, \emph{śikhar}.) A mound? (\textbf{\hi{टीला}}, \emph{ṭīlā}.) A heap? (\textbf{\hi{ढेर}}, \emph{ḍher}.) A straw? (\textbf{\hi{तिनका}}, \emph{tinkā}.) To roll (something)? (\textbf{\hi{लुढ़काना}}, \emph{luṛhkānā}.) A canal? (\textbf{\hi{नहर}}, \emph{nahar}.) A reservoir? (\textbf{\hi{जलाशय}}, \emph{jalāśay}.) A current? (\textbf{\hi{धारा}}, \emph{dhārā}.) To irrigate? (\textbf{\hi{सींचना}}, \emph{sī̃cnā}.) To weed? (\textbf{\hi{निराना}}, \emph{nirānā}.) A harbour? (\textbf{\hi{बंदरगाह}}, \emph{bandargāh}.) An anchor? (\textbf{\hi{लंगर}}, \emph{laṅgar}.) A bay? (\textbf{\hi{खाड़ी}}, \emph{khāṛī}.) A coast? (\textbf{\hi{तट}}, \emph{taṭ}.) A dip? (\textbf{\hi{डुबकी}}, \emph{ḍubkī}.)
\end{tcolorbox}
